\documentclass[11pt]{article}

\usepackage[margin=1in]{geometry}
\usepackage{amsmath,amssymb}
\usepackage{booktabs}
\usepackage{array}
\usepackage{enumitem}
\usepackage[hidelinks]{hyperref}

\setlist[itemize]{leftmargin=1.4em}
\setlist[enumerate]{leftmargin=1.6em}

\title{SYSC 5108 Exam Study\\Assignment 2, Question 4}
\author{Jesse Levine}
\date{}

\begin{document}
\maketitle

\section*{Question}

For the multivariate logistic regression model from Assignment 2, Question 4,
continuous input features are normalized to the range \([-1,1]\), missing
continuous values are imputed with the mean, and missing categorical values are
imputed with the mode. Use the trained model to predict whether each query
instance will make a repeat purchase.

\section*{Course and Textbook Cross-References}

\begin{itemize}
    \item \textbf{Chapter 3 slides, Logistic Regression Model:} for \(K=2\)
    classes, the model uses a linear score \(z=\mathbf{w}^{T}\mathbf{x}\) and
    converts it to a probability with the logistic sigmoid.
    \item \textbf{Chapter 3 slides, Logistic Sigmoid Function:}
    \[
        \sigma(z)=\frac{1}{1+\exp(-z)}.
    \]
    \item \textbf{Chapter 3 slides, Prediction:} the predicted class is selected
    by comparing class probabilities. With two classes, this is equivalent to
    predicting repeat purchase when \(\sigma(z)\ge 0.5\).
    \item \textbf{Goodfellow, Bengio, and Courville, \emph{Deep Learning},
    Chapter 5:} logistic regression maps a linear function through a sigmoid and
    interprets \(p(y=1\mid x;\theta)=\sigma(\theta^T x)\) as the probability of
    class 1.
\end{itemize}

\section*{Given Information}

\subsection*{Data Quality Report}

The continuous-feature statistics needed for range normalization and imputation
are:

\begin{center}
\begin{tabular}{lrrrrrr}
\toprule
Feature & Min & Mean & Max & 1st Qrt & Median & 3rd Qrt \\
\midrule
Age & 18.0 & 32.7 & 63.0 & 22.0 & 32.0 & 43.0 \\
Shop frequency & 0.2 & 2.2 & 5.4 & 1.0 & 1.3 & 4.3 \\
Shop value & 5.0 & 101.9 & 230.7 & 11.8 & 100.14 & 174.6 \\
\bottomrule
\end{tabular}
\end{center}

The categorical mode for socioeconomic band is
\[
    \operatorname{mode}(\text{socioeconomic band}) = a.
\]

\subsection*{Model Weights}

\begin{center}
\begin{tabular}{lr}
\toprule
Feature & Weight \\
\midrule
Intercept \(w_0\) & 0.6679 \\
Age & -0.5795 \\
Socioeconomic band B & -0.1981 \\
Socioeconomic band C & -0.2318 \\
Shop value & 3.4091 \\
Shop frequency & 2.0499 \\
\bottomrule
\end{tabular}
\end{center}

\subsection*{Query Instances}

\begin{center}
\begin{tabular}{rrrrr}
\toprule
ID & Age & Band & Shop frequency & Shop value \\
\midrule
1 & 38 & a & 1.90 & 165.39 \\
2 & 56 & b & 1.60 & 109.32 \\
3 & 18 & c & 6.00 & 10.09 \\
4 & ? & b & 1.33 & 204.62 \\
5 & 62 & ? & 0.85 & 110.50 \\
\bottomrule
\end{tabular}
\end{center}

\section*{Step-by-Step Solution}

\subsection*{1. Impute missing values}

Missing continuous values are imputed with the mean, so ID 4 has
\[
    \text{Age}_{4}=32.7.
\]
Missing categorical values are imputed with the mode, so ID 5 has
\[
    \text{Band}_{5}=a.
\]

\subsection*{2. Encode the categorical feature}

The model uses one-hot coding for bands \(b\) and \(c\). Band \(a\) is the baseline:
\[
    a \rightarrow (B,C)=(0,0),\qquad
    b \rightarrow (B,C)=(1,0),\qquad
    c \rightarrow (B,C)=(0,1).
\]

\subsection*{3. Normalize continuous features}

Range normalization to \([-1,1]\) is
\[
    x_{\text{norm}}
    =
    2\frac{x-x_{\min}}{x_{\max}-x_{\min}}-1.
\]

For example, for ID 1 age:
\[
    \text{Age}_{1,\text{norm}}
    =
    2\frac{38-18}{63-18}-1
    =
    -0.1111.
\]

\textbf{Assumption for ID 3:} the query value \(\text{Shop frequency}=6.00\)
is above the training-set max \(5.4\). Since the question says features are
normalized to \([-1,1]\), this solution clips ID 3's shop frequency to the
training max \(5.4\) before normalization. Without clipping, ID 3's normalized
frequency would be \(1.2308\) and its prediction changes; see the exam note at
the end.

\begin{center}
\begin{tabular}{rrrrrr}
\toprule
ID & Age norm & Freq. norm & Value norm & \(B\) & \(C\) \\
\midrule
1 & -0.1111 & -0.3462 & 0.4213 & 0 & 0 \\
2 & 0.6889 & -0.4615 & -0.0756 & 1 & 0 \\
3 & -1.0000 & 1.0000 & -0.9549 & 0 & 1 \\
4 & -0.3467 & -0.5654 & 0.7689 & 1 & 0 \\
5 & 0.9556 & -0.7500 & -0.0651 & 0 & 0 \\
\bottomrule
\end{tabular}
\end{center}

\subsection*{4. Compute the logistic-regression score}

Use
\[
\begin{aligned}
z
&=
w_0
+w_{\text{age}}x_{\text{age}}
+w_B x_B
+w_C x_C
+w_{\text{value}}x_{\text{value}}
+w_{\text{freq}}x_{\text{freq}}.
\end{aligned}
\]

For ID 1:
\[
\begin{aligned}
z_1
&=
0.6679
 +(-0.5795)(-0.1111)
 +(-0.1981)(0)
 +(-0.2318)(0) \\
&\quad
 +(3.4091)(0.4213)
 +(2.0499)(-0.3462) \\
&=1.4589.
\end{aligned}
\]

\subsection*{5. Convert the score to a probability}

Use the sigmoid:
\[
    p(\text{repeat purchase}\mid x)=\sigma(z)=\frac{1}{1+\exp(-z)}.
\]

Predict repeat purchase if \(p\ge 0.5\).

\begin{center}
\begin{tabular}{rrrr}
\toprule
ID & \(z\) & \(\sigma(z)\) & Predicted class \\
\midrule
1 & 1.4589 & 0.8114 & Repeat purchase \\
2 & -1.1332 & 0.2436 & No repeat purchase \\
3 & -0.1898 & 0.4527 & No repeat purchase \\
4 & 2.1330 & 0.8941 & Repeat purchase \\
5 & -1.6453 & 0.1617 & No repeat purchase \\
\bottomrule
\end{tabular}
\end{center}

\section*{Final Answer}

Using mean/mode imputation, one-hot encoding, range normalization, and the
clipped ID 3 shop-frequency value, the model predicts:
\[
\boxed{
\begin{array}{c|c|c}
\text{ID} & p(\text{repeat purchase}) & \text{Decision} \\
\hline
1 & 0.8114 & \text{repeat purchase} \\
2 & 0.2436 & \text{no repeat purchase} \\
3 & 0.4527 & \text{no repeat purchase} \\
4 & 0.8941 & \text{repeat purchase} \\
5 & 0.1617 & \text{no repeat purchase}
\end{array}}
\]

\section*{Exam Notes}

\begin{itemize}
    \item The most common mistakes in this question are forgetting to impute
    missing values before normalization, encoding band \(a\) incorrectly, or
    applying the sigmoid before adding the intercept.
    \item The score \(z\) is linear, but the final probability is nonlinear because
    of the sigmoid.
    \item If the instructor expects \emph{unclipped} range normalization for the
    out-of-range ID 3 frequency value \(6.00\), then
    \[
        x_{\text{freq},3}
        =
        2\frac{6.00-0.2}{5.4-0.2}-1
        =
        1.2308,
    \]
    which gives
    \[
        z_3=0.2832,\qquad \sigma(z_3)=0.5703.
    \]
    Under that convention, ID 3 would be predicted as repeat purchase. The
    clipped convention matches the assignment solution assumption that normalized
    features should remain inside \([-1,1]\).
\end{itemize}

\end{document}
